% AA-13 | Algebra abstrakcyjna | Iteracja VII
\chapter[Pierścienie noetherowskie]{Pierścienie noetherowskie i twierdzenie Hilberta o bazie}
\label{aa:ch:noether-hilbert}

\section{Dlaczego skończona liczba generatorów jest istotna?}
W pierścieniu $\mathbb Z$ każdy ideał jest wyznaczony przez
jeden element, natomiast w $K[x,y]$, gdzie $K$ jest ciałem,
nie każdy ideał jest główny.
Nie oznacza to jednak, że do opisania ideału potrzeba zawsze
nieskończenie wielu generatorów. Twierdzenie Hilberta o bazie
mówi, że przejście od pierścienia noetherowskiego do pierścienia
wielomianów jednej zmiennej zachowuje skończone generowanie
\emph{wszystkich} ideałów.

W rozdziale pracujemy z przemiennymi pierścieniami z jedynką,
jeśli nie wskazano inaczej. W przypadku pierścienia zerowego
stwierdzenia o skończonym generowaniu są natychmiastowe;
w przykładach geometrycznych przyjmujemy pierścienie niezerowe.
Przypomnijmy, że ideał
$I\triangleleft R$ jest skończenie generowany, jeśli
$I=(a_1,\ldots,a_r)$ dla pewnej skończonej rodziny elementów
$a_i\in R$. Dopuszczamy $r=0$; wówczas rodzina generatorów jest pusta,
a generowany przez nią ideał jest zerowy.

\section{Skończone generowanie a łańcuchy ideałów}
\begin{definicja}[Pierścień noetherowski]
\label{aa:def:noether-full}
Przemienny pierścień $R$ z jedynką jest \emph{noetherowski},
jeżeli każdy jego ideał jest skończenie generowany.
Równoważnie, na mocy
Twierdzenia~\ref{aa:tw:noetherian-acc},
każdy wstępujący ciąg ideałów
$I_1\subseteq I_2\subseteq\cdots$ stabilizuje się,
to znaczy istnieje $N$ takie, że $I_n=I_N$ dla każdego
$n\ge N$.
\end{definicja}

\begin{lemat}[Maksymalny element rodziny ideałów]
\label{aa:lem:noether-max-ideal-family}
Niech $R$ będzie pierścieniem noetherowskim, a
$\mathcal F$ niepustą rodziną ideałów $R$.
Wówczas $\mathcal F$ zawiera ideał maksymalny względem
zawierania w rodzinie $\mathcal F$.
\end{lemat}
\begin{proof}
Przypuśćmy, że żaden ideał $I_1\in\mathcal F$ nie jest
maksymalny. Wtedy istnieje $I_2\in\mathcal F$ ściśle
zawierający $I_1$. Skoro również $I_2$ nie jest maksymalny,
wybieramy $I_3\in\mathcal F$ z $I_2\subsetneq I_3$.
Powtarzając tę konstrukcję, otrzymujemy nieskończony
ściśle wstępujący ciąg ideałów. Przeczy to warunkowi
stabilizacji. Zatem pewien element $\mathcal F$ musi
być maksymalny.
\end{proof}

\begin{twierdzenie}[Ilorazy pierścieni noetherowskich]
\label{aa:tw:noether-quotients}
Jeśli $R$ jest pierścieniem noetherowskim i
$I\triangleleft R$ jest ideałem, to $R/I$ jest
pierścieniem noetherowskim.
\end{twierdzenie}
\begin{proof}
Niech $J$ będzie dowolnym ideałem $R/I$, a
$\pi:R\to R/I$ kanonicznym homomorfizmem ilorazowym.
Przeciwobraz $K=\pi^{-1}(J)$ jest ideałem $R$
zawierającym $I$. Z noetherowskości $R$ istnieją
$a_1,\ldots,a_r\in K$ takie, że $K=(a_1,\ldots,a_r)$.
Dla każdego $y\in J$ wybieramy $x\in K$ z $\pi(x)=y$
i zapisujemy $x=\sum_{i=1}^r t_ia_i$ z $t_i\in R$.
Po zastosowaniu $\pi$ otrzymujemy
$y=\sum_i\pi(t_i)\pi(a_i)$.
Zatem $J=(\pi(a_1),\ldots,\pi(a_r))$ w $R/I$.
Każdy ideał ilorazu ma skończenie wielu generatorów.
\end{proof}

\begin{twierdzenie}[Lokalizacje pierścieni noetherowskich]
\label{aa:tw:noether-localizations}
Niech $R$ będzie pierścieniem noetherowskim, a
$S\subseteq R$ zbiorem multiplikatywnym niezawierającym
zera. Wtedy lokalizacja $S^{-1}R$ jest noetherowska.
\end{twierdzenie}
\begin{proof}
Niech $J\triangleleft S^{-1}R$ i oznaczmy przez
$\iota:R\to S^{-1}R$ kanoniczny homomorfizm.
Przeciwobraz $I=\iota^{-1}(J)$ jest ideałem $R$,
więc $I=(a_1,\ldots,a_r)$ dla pewnych $a_i\in R$.
Wykażemy, że $J$ jest generowany przez ułamki $a_i/1$.
Każdy $a_i/1$ należy do $J$, więc zawieranie ideału
przez nie generowanego w $J$ jest oczywiste.
W drugą stronę, jeżeli $b/s\in J$, to
$b/1=(s/1)(b/s)\in J$, zatem $b\in I$.
Istnieją $c_1,\ldots,c_r\in R$ z
$b=\sum_i c_ia_i$. Dzieląc przez $s$ w lokalizacji,
otrzymujemy $b/s=\sum_i (c_i/s)(a_i/1)$.
Oznacza to, że $J=(a_1/1,\ldots,a_r/1)$.
\end{proof}

\begin{przyklad}[Pierścień lokalny uzyskany z wielomianów]
\label{aa:prz:noether-local-ring}
Niech $K$ będzie ciałem. Ponieważ $K[x,y]$ będzie noetherowski na mocy twierdzenia
Hilberta, jego lokalizacja
$K[x,y]_{(x,y)}$ także jest noetherowska.
Jest ona jednocześnie pierścieniem lokalnym: jedyny ideał
maksymalny $\mathfrak m=(x,y)K[x,y]_{(x,y)}$ jest generowany
przez obrazy $x$ i $y$.
Nie jest to ciało, ponieważ $x$ należy do $\mathfrak m$
i nie jest jednostką. Ponadto $\mathfrak m$ nie jest główny.
Istotnie, przestrzeń $\mathfrak m/\mathfrak m^2$ nad
$K\cong K[x,y]_{(x,y)}/\mathfrak m$ ma bazę złożoną z klas
$x$ i $y$: każdy licznik elementu $\mathfrak m$ ma
wyraz stały równy zeru, a mianownik ma niezerowy wyraz
stały. Dokładniej, jeśli $s=s_0+q$ z $s_0\in K^\times$
i $q\in(x,y)$, to
$s^{-1}-s_0^{-1}=-s_0^{-1}q/s\in\mathfrak m$.
Po pomnożeniu przez licznik należący do $\mathfrak m$
widzimy, że modulo $\mathfrak m^2$ ułamek jest kombinacją
liniową klas $x$ i $y$. Klasy są liniowo niezależne:
jeśli $ax+by\in\mathfrak m^2$ dla $a,b\in K$, to z
kryterium przynależności do rozszerzonego ideału
(Lemat~\ref{aa:lem:local-ideal-membership}) wynika,
że pewien wielomian $s$ z $s(0,0)\ne0$ spełnia
$s(ax+by)\in(x^2,xy,y^2)$. Porównanie składników
stopnia jeden daje $s(0,0)a=s(0,0)b=0$, a zatem
$a=b=0$. Zatem
$\dim_K(\mathfrak m/\mathfrak m^2)=2$. Gdyby
$\mathfrak m=(h)$, ten iloraz byłby generowany przez
jedną klasę $h$ nad $K$, więc miałby wymiar najwyżej
jeden, sprzeczność. Noetherowskość nie oznacza zatem
ani główności ideałów, ani odwracalności wszystkich
elementów niezerowych.
\end{przyklad}

\section{Moduły nad pierścieniami noetherowskimi}
Najpierw wyprowadzimy narzędzie dotyczące podmodułów
modułów skończenie generowanych. Następnie przedstawimy
bezpośredni dowód twierdzenia Hilberta o bazie oparty na
współczynnikach wiodących. Wyniki o podmodułach wykorzystamy
ponadto do badania algebr skończonych jako moduły.

\begin{definicja}[Moduł skończenie generowany]
\label{aa:def:noether-fg-module}
Niech $R$ będzie przemiennym pierścieniem z jedynką,
a $M$ modułem nad $R$. Mówimy, że $M$ jest
\emph{skończenie generowany}, jeśli dla pewnej liczby całkowitej
$r\ge0$ istnieją elementy $m_1,\ldots,m_r\in M$
takie, że każdy $m\in M$ jest
postaci $m=\sum_{i=1}^r a_im_i$ dla $a_i\in R$.
Gdy $r=0$, suma jest pusta i $M=0$.
\end{definicja}

\begin{lemat}[Podmoduły modułów wolnych skończonej rangi]
\label{aa:lem:noether-submodules-free}
Jeśli $R$ jest pierścieniem noetherowskim, to każdy
podmoduł $N$ modułu $R^r$ dla $r\ge0$ jest
skończenie generowany.
\end{lemat}
\begin{proof}
Indukujemy względem $r$. Dla $r=0$ moduł $R^0$ jest zerowy,
więc jedynym podmodułem jest moduł zerowy.
Dla $r=1$ podmoduły $R$ to dokładnie ideały, które są
skończenie generowane na mocy założenia.

Niech $r\ge2$ i niech $N\subseteq R^r$.
Rozważmy odwzorowanie $p:R^r\to R$, które wybiera
ostatnią współrzędną. Obraz $p(N)$ jest ideałem $R$,
więc ma generatory $a_1,\ldots,a_s$.
Dla każdego $a_i$ wybieramy $v_i\in N$ z
$p(v_i)=a_i$. Z drugiej strony
$N\cap(R^{r-1}\times\{0\})$ jest podmodułem
$R^{r-1}$ po pominięciu ostatniej współrzędnej,
więc na mocy założenia indukcyjnego ma generatory
$w_1,\ldots,w_t$. Dla dowolnego $v\in N$ zapisujemy
$p(v)=\sum_i c_i a_i$. Wtedy
$v-\sum_i c_i v_i\in N$ ma ostatnią współrzędną zerową,
więc jest kombinacją $w_j$. W rezultacie elementy
$v_i,w_j$ generują cały $N$.
\end{proof}

\begin{twierdzenie}[Podmoduły modułów skończenie generowanych]
\label{aa:tw:noether-submodules-fg}
Niech $R$ będzie przemiennym pierścieniem noetherowskim
z jedynką. Jeśli $M$ jest skończenie generowanym
$R$-modułem, to każdy podmoduł $N\subseteq M$
jest skończenie generowany.
\end{twierdzenie}
\begin{proof}
Wybierzmy generatory $m_1,\ldots,m_r$ modułu $M$.
Homomorfizm $\varphi:R^r\to M$ dany przez
$\varphi(a_1,\ldots,a_r)=\sum_i a_im_i$ jest
surjektywny. Przeciwobraz $\varphi^{-1}(N)$ jest
podmodułem $R^r$, więc poprzedni lemat daje jego
skończony układ generatorów $u_1,\ldots,u_s$.
Ich obrazy $\varphi(u_i)$ generują $N$: dla każdego
$n\in N$ wybieramy $u\in R^r$ z $\varphi(u)=n$,
rozkładamy $u$ na sumę kombinacji $u_i$ i stosujemy
$\varphi$. Tak otrzymujemy wymagane przedstawienie.
\end{proof}

\begin{przyklad}[Dlaczego wymagana jest noetherowskość]
\label{aa:prz:noether-need}
Niech $K$ będzie ciałem i niech $R=K[x_1,x_2,\ldots]$ będzie pierścieniem
wielomianów w przeliczalnie wielu zmiennych.
Jako moduł nad sobą $R$ jest generowany przez $1$.
Podmoduł $I=(x_1,x_2,\ldots)$ nie jest jednak
skończenie generowany. Gdyby generowały go
$f_1,\ldots,f_r$, to wszystkie te wielomiany
występowałyby jedynie w skończonej liczbie zmiennych
$x_1,\ldots,x_N$. Po podstawieniu tych zmiennych
równych zeru obrazy generatorów byłyby zerowe
(każdy $f_i$ ma wyraz wolny równy zeru), podczas gdy
obraz $x_{N+1}$ pozostałby niezerowy. Sprzeczność.
\end{przyklad}


\begin{twierdzenie}[Skończone prezentacje modułów]
\label{aa:tw:noether-finite-presentation}
Niech $R$ będzie przemiennym pierścieniem noetherowskim
z jedynką, a $M$ skończenie generowanym $R$-modułem.
Wtedy istnieją liczby całkowite $r,s\geq0$ oraz
homomorfizmy $R$-modułów
$R^s\xrightarrow{\,A\,}R^r\xrightarrow{\,\varphi\,}M$
takie, że $\varphi$ jest surjekcją oraz
$\operatorname{im} A=\ker\varphi$.
Innymi słowy, moduł $M$ ma skończony układ generatorów
i skończony układ relacji między nimi.
\end{twierdzenie}
\begin{proof}
Wybierzmy generatory $m_1,\ldots,m_r$ modułu $M$.
Zdefiniujmy $\varphi:R^r\to M$ wzorem
$\varphi(c_1,\ldots,c_r)=\sum_{i=1}^r c_i m_i$.
Definicja generowania gwarantuje, że $\varphi$ jest
surjekcją. Jądro $K=\ker\varphi$ jest podmodułem
$R^r$, a z Lematu~\ref{aa:lem:noether-submodules-free}
wynika, że istnieją jego generatory $v_1,\ldots,v_s$.
Niech $A:R^s\to R^r$ odwzorowuje $j$-ty standardowy
wektor bazowy na $v_j$. Wówczas obraz $A$ jest
podmodułem generowanym przez $v_j$, czyli dokładnie
$K$. Otrzymaliśmy $\operatorname{im}A=\ker\varphi$,
co jest żądaną skończoną prezentacją.
Jeśli $K=0$, można przyjąć $s=0$.
\end{proof}

\begin{przyklad}[Prezentacja modułu $\mathbb Z/6\mathbb Z$]
\label{aa:prz:noether-finite-presentation-z6}
Dla pierścienia $R=\mathbb Z$ i $M=\mathbb Z/6\mathbb Z$
bierzemy $\varphi:\mathbb Z\to M$, $n\mapsto\overline n$.
Jego jądrem jest $6\mathbb Z$, więc
$A:\mathbb Z\to\mathbb Z$, $A(n)=6n$, ma obraz
równy $\ker\varphi$. Otrzymujemy ciąg dokładny
$\mathbb Z\xrightarrow{\,\cdot6\,}\mathbb Z
\xrightarrow{\,\varphi\,}\mathbb Z/6\mathbb Z\to0$.
Jeden generator $\overline1$ wystarcza, lecz spełnia
nietrywialną relację $6\overline1=0$; z tego powodu
moduł nie jest wolny rangi jeden.
\end{przyklad}

\section{Twierdzenie Hilberta o bazie}
\begin{twierdzenie}[Twierdzenie Hilberta o bazie]
\label{aa:tw:hilbert-basis}
Jeżeli $R$ jest przemiennym pierścieniem noetherowskim
z jedynką, to pierścień wielomianów $R[x]$ jest
noetherowski.
\end{twierdzenie}
\begin{proof}
Niech $I\triangleleft R[x]$ będzie dowolnym ideałem.
Jeżeli $I=(0)$, to jest generowany przez pustą rodzinę.
Załóżmy więc $I\ne(0)$. Zdefiniujmy $L\subseteq R$
jako zbiór zawierający $0$ oraz wszystkie współczynniki
wiodące niezerowych wielomianów należących do $I$.

Pokażemy, że $L$ jest ideałem $R$. Jeśli $a,b\in L$
i nie są zerami, wybierzmy $f,g\in I$ o współczynnikach
wiodących odpowiednio $a,b$ i stopniach $m,n$.
Mnożąc $f$ przez $x^{\max(m,n)-m}$, a $g$ przez
$x^{\max(m,n)-n}$, otrzymujemy wielomiany
tego samego stopnia, których suma ma
współczynnik wiodący $a+b$, gdy $a+b\ne0$.
W przypadku $a+b=0$ wynik należy do $L$ z definicji.
Podobnie dla $r\in R$ współczynnik wiodący $rf$
równa się $ra$, jeśli $ra\ne0$, a w przeciwnym razie
$ra=0\in L$. Również elementy przeciwne należą do $L$.
Zatem $L$ jest ideałem.

Ponieważ $R$ jest noetherowski, mamy
$L=(a_1,\ldots,a_s)$. Wybieramy dla każdego
niezerowego generatora $a_i$ wielomian $f_i\in I$
o współczynniku wiodącym $a_i$ i stopniu $d_i$.
Generatory zerowe można pominąć. Niech
$D=\max(d_1,\ldots,d_s)$. Ponieważ $I\ne(0)$,
ideał $L$ jest niezerowy i lista wybranych wielomianów
nie jest pusta.

Dla każdego $j$ z $0\le j<D$ definiujemy
$L_j\subseteq R$ jako zbiór zawierający $0$ oraz
współczynniki wiodące wielomianów $h\in I$ o
stopniu dokładnie $j$. Powyższy argument bez
przesuwania stopni pokazuje, że każdy $L_j$ jest
ideałem (gdy dwa współczynniki sumują się do zera,
zero już należy do zbioru). Wybieramy skończony układ
generatorów $b_{j,1},\ldots,b_{j,r_j}$ ideału $L_j$
i wielomiany $g_{j,k}\in I$ stopnia $j$ o odpowiednich
współczynnikach wiodących. Jeśli $L_j=0$, nie wybieramy
żadnego wielomianu.

Udowodnimy, że skończona rodzina wszystkich $f_i$ i
$g_{j,k}$ generuje $I$. Weźmy niezerowy wielomian
$h\in I$ o stopniu $n$ i współczynniku wiodącym $c$.
Jeżeli $n\ge D$, piszemy $c=\sum_i r_i a_i$,
a następnie odejmujemy od $h$ wielomian
$\sum_i r_i x^{n-d_i}f_i\in I$. Wyrazy stopnia $n$
się redukują, więc pozostały wielomian jest zerowy
albo ma stopień mniejszy od $n$.
Jeśli natomiast $n<D$, to $c\in L_n$, zatem
$c=\sum_k r_k b_{n,k}$. Odejmując
$\sum_k r_k g_{n,k}$, ponownie obniżamy stopień.

W obu przypadkach uzyskujemy wielomian w $I$
o ściśle mniejszym stopniu. Powtarzamy krok aż do
zera, co musi nastąpić, bo stopnie niezerowych
wielomianów są nieujemne i nie mogą maleć bez końca.
Odjęte wielomiany należą do ideału generowanego przez
skończoną rodzinę $f_i,g_{j,k}$; zatem także
początkowy $h$ należy do tego ideału. Dowodzi to,
że $I$ jest skończenie generowany. Ponieważ $I$
był dowolny, $R[x]$ jest noetherowski.
\end{proof}

\begin{wniosek}[Pierścienie wielomianów skończenie wielu zmiennych]
\label{aa:wn:hilbert-finite-vars}
Jeśli $R$ jest pierścieniem noetherowskim, to
$R[x_1,\ldots,x_n]$ jest noetherowski dla każdego
$n\ge0$. W szczególności $K[x_1,\ldots,x_n]$
jest noetherowski dla każdego ciała $K$.
\end{wniosek}
\begin{proof}
Dla $n=0$ otrzymujemy sam pierścień $R$.
Gdy wynik jest prawdziwy dla $n-1$, zastosowanie
twierdzenia Hilberta do
$R[x_1,\ldots,x_{n-1}]$ daje noetherowskość
$(R[x_1,\ldots,x_{n-1}])[x_n]$, czyli
$R[x_1,\ldots,x_n]$. Ciało jest noetherowskie,
ponieważ jego jedynymi ideałami są $(0)$ i całe ciało.
\end{proof}

\begin{przyklad}[Skończone generowanie nie oznacza jednego generatora]
\label{aa:prz:hilbert-not-pid}
Niech $K$ będzie ciałem. Ideał $(x,y)\triangleleft K[x,y]$ ma dwa generatory.
Nie ma jednego generatora, co wykazaliśmy w
Przykładzie~\ref{aa:prz:two-polynomials-not-pid}.
Twierdzenie Hilberta gwarantuje jedynie, że dowolny
ideał $K[x,y]$ posiada \emph{jakąś skończoną}
rodzinę generatorów. Nie gwarantuje, że wystarczy jeden.
\end{przyklad}


\begin{twierdzenie}[Noetherowskość algebr skończonych jako moduły]
\label{aa:tw:noether-module-finite-algebra}
Niech $R$ będzie przemiennym pierścieniem noetherowskim
z jedynką, a $S$ przemiennym pierścieniem z jedynką,
wyposażonym w homomorfizm pierścieni
$\eta:R\to S$ zachowujący jedynkę. Jeśli $S$ jest
skończenie generowanym $R$-modułem dla działania
$r\cdot s=\eta(r)s$, to $S$ jest pierścieniem
noetherowskim.
\end{twierdzenie}
\begin{proof}
Niech $J$ będzie dowolnym ideałem pierścienia $S$.
Ponieważ $\eta(r)s\in J$ dla $r\in R$ i $s\in J$,
ideał $J$ jest też $R$-podmodułem $S$.
Twierdzenie~\ref{aa:tw:noether-submodules-fg}
stosuje się do $R$-modułu $S$, ponieważ jest on
skończenie generowany. Otrzymujemy
$J=Ra_1+\cdots+Ra_m$ dla pewnych $a_i\in J$.
Każde $a_i$ należy do $J$, zatem ideał $S$
generowany przez te elementy jest zawarty w $J$.
Z drugiej strony
$J=\sum_i\eta(R)a_i\subseteq\sum_i Sa_i$.
Stąd $J=(a_1,\ldots,a_m)$ jako ideał $S$.
Ponieważ $J$ był dowolny, $S$ jest noetherowski.
Nie wymagamy, by $\eta$ był iniektywny.
\end{proof}

\begin{przyklad}[Podalgebra generowana przez $t^2$ i $t^3$]
\label{aa:prz:noether-cusp-algebra}
Niech $K$ będzie ciałem, $R=K[t^2]$ i
$S=K[t^2,t^3]\subseteq K[t]$. Pierścień $R$
jest izomorficzny z $K[u]$ przez odwzorowanie
$u\mapsto t^2$, zatem jest noetherowski.
Pokażemy, że $S=R+Rt^3$ jako $R$-moduł.
Każdy iloczyn potęg generatorów $t^2,t^3$ ma postać
$t^{2a+3b}$ dla $a,b\ge0$. Gdy $b$ jest parzyste,
$t^{2a+3b}\in R$; gdy $b$ jest nieparzyste,
$t^{2a+3b}=t^3(t^2)^{a+3(b-1)/2}\in Rt^3$.
W drugą stronę zarówno $R$, jak i $Rt^3$ są
podzbiorami $S$. Stąd $S=R+Rt^3$ i twierdzenie
daje noetherowskość $S$.
Jednocześnie $t\notin S$: żaden z jednomianów
generujących $S$ nie ma stopnia jeden, więc
wielomian $t$ nie może być ich kombinacją liniową.
To pokazuje, że skończone generowanie $S$ nad $R$
nie wymaga równości $S=K[t]$.
\end{przyklad}

\begin{przyklad}[Dlaczego same współczynniki wiodące nie wystarczają]
\label{aa:prz:noether-leading-coefficients-example}
Rozważmy ideał $I=(2,x)$ w $\mathbb Z[x]$.
Wszystkie współczynniki wiodące niezerowych
wielomianów z $I$ generują ideał $L=\mathbb Z$:
ponieważ $x\in I$ ma współczynnik wiodący równy $1$,
wystarczy jeden generator $1$.
Wybrany dla niego wielomian $f=x$ ma stopień $D=1$,
lecz sam $x$ nie generuje całego $I$, ponieważ
$2\notin(x)$.
W konstrukcji z dowodu twierdzenia Hilberta potrzeba
jeszcze ideału współczynników stałych $L_0=(2)$,
któremu odpowiada wielomian $g=2$.
Rzeczywiście, jeśli $h=c_0+c_1x+\cdots+c_nx^n\in I$,
to $c_0=h(0)$ jest parzyste, gdyż $h=2a+xb$
dla pewnych $a,b\in\mathbb Z[x]$. Odwrotnie,
gdy $c_0=2k$, wielomian $q=(h-c_0)/x$ należy
do $\mathbb Z[x]$ (także dla $h$ stałego, gdy $q=0$),
a więc $h=2k+xq\in(2,x)$.
Przykładowo $3x^2+5x+6=x(3x+5)+3\cdot2$.
To obliczenie wyjaśnia rolę osobnych generatorów
$L_j$ dla stopni mniejszych niż $D$.
\end{przyklad}

\section{Ideały pierwotne: pierwsze przykłady}
\begin{definicja}[Ideał pierwotny]
\label{aa:def:primary-ideal}
Niech $R$ będzie przemiennym pierścieniem z jedynką.
Właściwy ideał $Q\subsetneq R$ nazywamy
\emph{pierwotnym}, jeśli dla wszystkich $a,b\in R$
z warunków $ab\in Q$ oraz $a\notin Q$ wynika,
że $b^n\in Q$ dla pewnego $n\ge1$.
\end{definicja}

\begin{lemat}[Pierwiastnik ideału pierwotnego]
\label{aa:lem:primary-radical-prime}
Niech $R$ będzie przemiennym pierścieniem z jedynką,
a $Q\subsetneq R$ ideałem pierwotnym. Wówczas $\sqrt Q$
jest ideałem pierwszym pierścienia $R$.
\end{lemat}
\begin{proof}
Pierwiastnik $\sqrt Q$ jest ideałem na mocy
Twierdzenia~\ref{aa:tw:radical-ideal}; jest właściwy,
bo $1\in\sqrt Q$ oznaczałoby $1\in Q$.
Niech $ab\in\sqrt Q$. Wtedy $(ab)^n\in Q$
dla pewnego $n\ge1$. Jeśli $a\in\sqrt Q$, teza
o pierwszości jest spełniona. Załóżmy więc
$a\notin\sqrt Q$, czyli $a^n\notin Q$.
Ponieważ $a^n b^n\in Q$ i $Q$ jest pierwotny,
istnieje $k\ge1$ takie, że $(b^n)^k\in Q$.
Stąd $b\in\sqrt Q$. To dowodzi pierwszości.
\end{proof}

\begin{przyklad}[Rozkład ideału na ideały pierwotne]
\label{aa:prz:primary-decomposition-explicit}
W pierścieniu $K[x,y]$ nad ciałem $K$ zachodzi
\[
 (x^2,xy)=(x)\cap(x^2,y).
\]
Uzasadnijmy obie inkluzje. Generatory $x^2$ i $xy$
należą do $(x)$ oraz $(x^2,y)$, więc lewa strona
jest zawarta w przecięciu. W drugą stronę, jeśli
$f\in(x)\cap(x^2,y)$, to $f=xg$ i
$f=x^2 a+yb$ dla pewnych $g,a,b\in K[x,y]$.
Po podstawieniu $x=0$ otrzymujemy
$0=y\,b(0,y)$ w dziedzinie $K[y]$, więc $b(0,y)=0$.
Oznacza to, że $x$ dzieli $b$, czyli $b=xc$.
Wówczas $f=x^2 a+xyc\in(x^2,xy)$.

Ideał $(x)$ jest pierwszy, ponieważ
$K[x,y]/(x)\cong K[y]$ jest dziedziną całkowitości.
Ideał $(x^2,y)$ jest pierwotny. Istotnie, jego iloraz
jest izomorficzny z $K[x]/(x^2)$. W tym ostatnim
pierścieniu każda klasa ma jednoznaczną postać
$a+bx+(x^2)$ z $a,b\in K$. Jeśli $a\ne0$, klasa
jest jednostką: jej odwrotność wynosi
$a^{-1}-a^{-2}bx+(x^2)$. Jeśli $a=0$, klasa ma
kwadrat równy zeru. Jeśli iloczyn klas $u,v$ jest zerowy,
a $u$ jest niezerowa, to $v$ nie może być jednostką,
bo pomnożenie przez $v^{-1}$ dałoby $u=0$.
Zatem $v^2=0$, co jest warunkiem pierwotności
ideału $(x^2,y)$.

Jest to \emph{konkretny przykład} rozkładu na ideały
pierwotne, a nie jeszcze dowód ogólnego twierdzenia
o rozkładzie pierwotnym ideałów pierścieni noetherowskich.
Pełną teorię tej dekompozycji rozwiniemy dopiero po
ustaleniu dalszych narzędzi modułowych.
\end{przyklad}
